\section{Desain dan Konstruksi AST}
\label{sec:ch16-desain}

\subsection{Prinsip Desain}

AST yang baik memenuhi beberapa prinsip:
\begin{itemize}[leftmargin=*]
  \item setiap simpul merepresentasikan satu konstruksi bahasa;
  \item informasi sintaksis yang tidak bermakna dihilangkan;
  \item simpul menyimpan lokasi untuk diagnostik;
  \item struktur mendukung traversal yang mudah;
  \item tipe dan informasi semantik dapat ditambahkan tanpa mengubah
        bentuk dasar.
\end{itemize}

\subsection{Struktur Generik}

Proyek Subset-C memakai simpul generik:
\begin{lstlisting}[language=C]
typedef struct AstNode {
    AstKind kind;
    Token token;
    TypeKind decl_type;
    struct AstNode *left;
    struct AstNode *right;
    struct AstNode *children;
    struct AstNode *next;
} AstNode;
\end{lstlisting}

Struktur ini menyimpan:
\begin{itemize}[leftmargin=*]
  \item \code{kind}: jenis simpul;
  \item \code{token}: token representatif untuk lokasi dan lexeme;
  \item \code{decl\_type}: tipe deklarasi (untuk deklarasi/fungsi);
  \item \code{left} dan \code{right}: operand biner/unary;
  \item \code{children} dan \code{next}: daftar anak.
\end{itemize}

Antarmuka lengkap AST ditunjukkan pada \lstref{lst:ch16-ast-h}.

\includecode[
  language=C,
  firstline=1,
  lastline=70,
  firstnumber=1,
  caption={Antarmuka \file{ast.h}},
  label={lst:ch16-ast-h}
]{../../code/subsetc/include/ast.h}

\subsection{Generik versus Spesifik}

Simpul generik lebih mudah diimplementasikan dan dipelajari, tetapi
kurang type-safe. Simpul spesifik, misalnya \code{BinaryExpr} dan
\code{IfStmt}, memberikan kejelasan dan mencegah akses field yang
salah. Untuk proyek besar, simpul spesifik biasanya lebih baik.

\begin{infobox}
Pilihan struktur AST memengaruhi seluruh fase berikutnya. Luangkan
waktu untuk merancangnya sebelum menulis parser. Perubahan struktur
setelah parser selesai biasanya mahal.
\end{infobox}

\subsection{Membangun AST}

AST dibangun saat parsing:
\begin{itemize}[leftmargin=*]
  \item aturan yang hanya mengatur sintaksis tidak membuat simpul;
  \item aturan operator membuat simpul biner dengan dua anak;
  \item aturan daftar membuat simpul dan menambahkan anak;
  \item deklarasi menyimpan tipe pada \code{decl\_type}.
\end{itemize}

Pendekatan ini menghasilkan AST yang ringkas dan langsung mencerminkan
struktur program.
